%! Rational Functions and End Behavior — Unit 3, Lesson 3.1 — Slides (Beamer)
\documentclass[aspectratio=169, 11pt]{beamer}
\usepackage{precalculus-beamer}   % provides \navyheader{} and \sectionlabel[color]{}

\begin{document}

% TITLE SLIDE
{
\setbeamercolor{background canvas}{bg=navy}
\begin{frame}[plain]
  \vspace{2.8cm}\hspace{0.55cm}%
  \begin{minipage}{0.9\textwidth}
    {\color{goldacc}\bfseries\small LESSON 1.7}\\[0.5cm]
    {\color{white}\bfseries\LARGE Rational Functions and End Behavior}\\[1.2cm]
    {\color{skymid}\small Precalculus~$\cdot$~Unit 3: Rational Functions}
  \end{minipage}
\end{frame}
}

% HOOK SLIDE
\begin{frame}[t, plain]
  \navyheader{Hook: Therapy Costs}
  \sectionlabel[goldacc]{Think About It}

  A therapist charges \$200 per session plus a one-time \$50 intake fee.\\
  \textbf{Average cost per session:}
  \[
    r(n) = \frac{200n + 50}{n}
  \]

  \bigskip
  \begin{columns}[T]
    \column{0.46\textwidth}
      \begin{block}{Compute some values}
        \begin{itemize}\setlength{\itemsep}{5pt}
          \item $r(1) = \$250$
          \item $r(10) = \$205$
          \item $r(100) = \$200.50$
          \item $r(1000) = \$200.05$
        \end{itemize}
      \end{block}

    \column{0.48\textwidth}
      \begin{block}{Discussion}
        \begin{itemize}\setlength{\itemsep}{5pt}
          \item What value does the average cost approach?
          \item Can it ever equal exactly \$200?
          \item What is the quotient of the leading terms?
        \end{itemize}
      \end{block}
  \end{columns}

  \bigskip
  \textit{Takeaway:} The graph has a \textbf{horizontal asymptote at $y = 200$}
  because $200n/n = 200$ as $n \to \infty$.
\end{frame}

% QUOTIENT OF LEADING TERMS
\begin{frame}[t, plain]
  \navyheader{Quotient of Leading Terms}
  \sectionlabel[navy]{Worked Example}

  For large $|x|$, a polynomial is dominated by its \textbf{leading term}.

  \medskip
  \textbf{Example:} $r(x) = \dfrac{3x^2 + 1}{x^2 - 4}$

  \bigskip
  \begin{columns}[T]
    \column{0.44\textwidth}
      \begin{block}{Step-by-Step}
        \begin{enumerate}\setlength{\itemsep}{5pt}
          \item Leading term of num: $3x^2$
          \item Leading term of den: $x^2$
          \item Quotient: $\dfrac{3x^2}{x^2} = 3$
          \item Horizontal asymptote: $y = 3$
        \end{enumerate}
      \end{block}

    \column{0.50\textwidth}
      \begin{block}{In Limit Notation}
        \[
          \lim_{x\to\infty} r(x) = 3
        \]
        \[
          \lim_{x\to -\infty} r(x) = 3
        \]
        \smallskip
        Check: $r(100) = \dfrac{30{,}001}{9{,}996} \approx 3.001$
      \end{block}
  \end{columns}
\end{frame}

% THE THREE CASES
\begin{frame}[t, plain]
  \navyheader{The Three End-Behavior Cases}

  Compare the \textbf{degree of the numerator} to the \textbf{degree of the denominator}:

  \bigskip
  \begin{block}{Summary Table}
    \renewcommand{\arraystretch}{1.6}
    \begin{tabular}{@{}lll@{}}
      \textbf{Degree Comparison} & \textbf{End Behavior} & \textbf{Example} \\
      \hline
      deg num $<$ deg denom & HA at $y = 0$ & $\dfrac{x+1}{x^2-9}$ \\[4pt]
      deg num $=$ deg denom & HA at $y = a_n/b_n$ & $\dfrac{3x^2}{x^2+4}$ \\[4pt]
      deg num $>$ deg denom & Polynomial end behavior & $\dfrac{x^3-2x}{x-3}$ \\[4pt]
    \end{tabular}
  \end{block}

  \bigskip
  \textbf{Special case:} If deg num $=$ deg denom $+ 1$, the graph has a
  \textbf{slant asymptote} (a non-horizontal line).
\end{frame}

% LIMIT NOTATION FOR HAs
\begin{frame}[t, plain]
  \navyheader{Limit Notation for Horizontal Asymptotes}
  \sectionlabel[navy]{EK 1.7.A.6}

  When a rational function $r$ has a horizontal asymptote at $y = b$:
  \[
    \lim_{x\to\infty} r(x) = b \qquad \text{and} \qquad \lim_{x\to -\infty} r(x) = b
  \]

  \bigskip
  \begin{columns}[T]
    \column{0.46\textwidth}
      \begin{block}{Equal Degrees}
        $r(x) = \dfrac{3x^2 + 1}{x^2 - 4}$\\[6pt]
        $\displaystyle\lim_{x\to\pm\infty} r(x) = 3$\\[4pt]
        HA: $y = 3$
      \end{block}

    \column{0.46\textwidth}
      \begin{block}{Num Degree $<$ Den Degree}
        $r(x) = \dfrac{x+1}{x^2-9}$\\[6pt]
        $\displaystyle\lim_{x\to\pm\infty} r(x) = 0$\\[4pt]
        HA: $y = 0$
      \end{block}
  \end{columns}
\end{frame}

% PAIR PRACTICE
\begin{frame}[t, plain]
  \navyheader{Class Practice}
  \sectionlabel[redacc]{Try It}

  With your partner, analyze each rational function.
  For each: identify the quotient of leading terms, state the asymptote type,
  and write both limit statements.

  \bigskip
  \begin{columns}[T]
    \column{0.46\textwidth}
      \begin{block}{Function 1}
        \[
          f(x) = \frac{4x^3 - 2}{x^3 + 5x}
        \]
        \begin{itemize}\setlength{\itemsep}{4pt}
          \item Quotient of leading terms: ?
          \item Asymptote type: ?
          \item $\lim_{x\to\infty} f(x) = $?
        \end{itemize}
      \end{block}

    \column{0.46\textwidth}
      \begin{block}{Function 2}
        \[
          g(x) = \frac{x^2 - 3}{x^4 + 1}
        \]
        \begin{itemize}\setlength{\itemsep}{4pt}
          \item Quotient of leading terms: ?
          \item Asymptote type: ?
          \item $\lim_{x\to -\infty} g(x) = $?
        \end{itemize}
      \end{block}
  \end{columns}
\end{frame}

\end{document}
